\section{LLM 数据墙与机器人数据荒漠的差异}

同样说“缺数据”，LLM 和机器人缺的不是同一种东西。LLM 的问题更像高质量公开语料边际收益下降：网页、书籍、代码、论文和问答数据已经被大规模使用，继续扩展会遇到重复、污染、版权、质量和评估瓶颈。机器人则更像从一开始就缺真实交互数据：没有足够多机器人在真实环境中执行任务，就没有稳定的端侧数据流。

\subsection{两类数据短缺}

\noindent\begin{longtable}{p{0.22\textwidth}p{0.35\textwidth}p{0.34\textwidth}}
\toprule
对象 & 数据短缺类型 & 典型解决路径 \\
\midrule
LLM & 高质量语料边际收益下降，benchmark 变难 & 更好过滤、更好合成数据、更高阶反馈、agent 环境。 \\
机器人 & 真实动作轨迹少、长尾场景少、采集昂贵 & 仿真、遥操作、小规模真实部署、失败数据挖掘。 \\
世界模型 & 物理预测数据和空间动态不足 & 视频、3D、仿真、多模态预测任务。 \\
VLA & 指令、视觉、动作、反馈难统一 & 端侧轨迹、仿真任务、语言标注、action head。 \\
\bottomrule
\end{longtable}

\begin{importantbox}{“数据墙”不是一个统一概念}
LLM 的数据墙是边际收益和质量问题；机器人的数据荒漠是物理交互和部署规模问题。混用这两个概念，会导致错误的数据策略。
\end{importantbox}

\subsection{本章小结}

LLM 和机器人都缺数据，但缺法不同。前者需要更高质量、更高阶反馈；后者需要可规模化的物理交互和仿真闭环。

